\section{A Torelli consequence}
\label{sec:torelli}
Hodge conservation applies to every member of the nine families.
Recovering a hypersurface from the resulting rational Hodge structure uses
a separate generic Torelli theorem, and its very-general qualification.

\begin{corollary}[Very-general cancellation]
\label{cor:generic-cancellation}\coverage{absent}%
Let $X\subset\PP^4$ be a very general smooth cubic or quartic threefold,
and let $Y\subset\PP^4$ be any smooth hypersurface of the same degree.
Then
\[
 X\times\PP^1\text{ is birational to }Y\times\PP^1
 \quad\Longleftrightarrow\quad X\cong Y.
\]
\end{corollary}
\begin{proof}
Theorem~\ref{thm:hodge-conservation} identifies the rational third Hodge
structures. Here $H^3$ is the entire primitive middle cohomology.
Voisin's rational generic Torelli theorem applies to both degrees
\cite[Theorem~0.2 and Remarks~0.1, 0.3]{VoisinTorelli}, and gives $X\cong Y$.
The reverse implication is immediate.
\uses{thm:hodge-conservation}%
\imports{Voisin:rational-torelli}%
\end{proof}

The arithmetic consequence in Appendix~\ref{sec:arithmetic-partners} uses
Hodge conservation with bounded-degree isogeny and polarization-finiteness
theorems. Neither application strengthens the integral or polarized content
of Theorem~\ref{thm:hodge-conservation}.
